% models/mathematical/planetare_zyklen_fourier.tex
% Fourier-Analyse planetarer Zyklen im Zeitfeldmodell

\section{Fourier-Analyse planetarer Zyklen}

\subsection{Grundidee}

Planetare Zyklen lassen sich als Summe harmonischer Komponenten darstellen:
\[
Z(t) = \sum_{k=1}^{N} A_k \cos(\omega_k t + \phi_k).
\]

Diese Darstellung bildet die Grundlage für die planetare Schicht des Zeitfeldes.

\subsection{Synodische Perioden}

Die synodische Frequenz zweier Planeten $i$ und $j$ ergibt sich aus:
\[
\omega_{ij} = |\omega_i - \omega_j|.
\]

Die zugehörige Fourier-Komponente:
\[
Z_{ij}(t) = A_{ij} \cos(\omega_{ij} t + \phi_{ij}).
\]

\subsection{Milanković-Komponenten}

Die drei Hauptzyklen:

\begin{align*}
\text{Exzentrizität:} &\quad \omega_E \approx \frac{2\pi}{100\,000 \text{ a}}, \\
\text{Obliquität:} &\quad \omega_O \approx \frac{2\pi}{41\,000 \text{ a}}, \\
\text{Präzession:} &\quad \omega_P \approx \frac{2\pi}{23\,000 \text{ a}}.
\end{align*}

Die Fourier-Darstellung:
\[
Z_{\text{Mil}}(t) = A_E \cos(\omega_E t) +
A_O \cos(\omega_O t) +
A_P \cos(\omega_P t).
\]

\subsection{Interferenz planetarer Frequenzen}

Die Interferenz zweier planetarer Zyklen:
\[
Z_{\text{int}}(t) = Z_i(t) + Z_j(t)
= A_i \cos(\omega_i t) + A_j \cos(\omega_j t).
\]

Dies ergibt Schwebungen:
\[
Z_{\text{beat}}(t) = 2 A \cos\left(\frac{\omega_i - \omega_j}{2} t\right)
\cos\left(\frac{\omega_i + \omega_j}{2} t\right).
\]

\subsection{Planetare Kopplung im Zeitfeld}

Die planetare Schicht des Zeitfeldtensors:
\[
\mathcal{T}^{(p)}_{ij}(t) = \sum_k A_k \cos(\omega_k t + \phi_k).
\]

\subsection{Fourier-Spektrum}

Das Spektrum ergibt sich aus:
\[
\tilde{Z}(\omega) = \int_{-\infty}^{\infty} Z(t) e^{-i\omega t} dt.
\]

Die Peaks entsprechen den planetaren Frequenzen.

\subsection{Kopplung an solare und galaktische Felder}

Die Kopplung erfolgt über:
\[
\mathcal{K}\big(\mathcal{T}^{(p)}, \mathcal{T}^{(s)}\big)
= \lambda_{ps} \sum_{k,l} A_k B_l \cos(\omega_k t + \phi_k)
\cos(\Omega_l t + \psi_l).
\]

\subsection{Ergebnis}

Die Fourier-Analyse bildet die mathematische Grundlage der planetaren
Zeitfeldkomponenten und ihrer Interferenzstrukturen.
